\documentclass[11pt]{article}
\usepackage{macrostyle}

\begin{document}

\lessonheader{Lectures 1--3: The Solow Growth Model}{Setup, Dynamics, Evaluation, Differential Equations, and Policy}{ECON*6020, Fall 2026}

\begin{recap}[How to read these notes]
The notes follow the four slide decks in order. Each section starts with a grey line telling you which slides it covers:
\textbf{L1} = \emph{The Solow Growth Model} (15 slides),
\textbf{EV} = \emph{Evaluating the Solow Growth Model} (18),
\textbf{DE} = \emph{Differential equations / speed of convergence} (14),
\textbf{PO} = \emph{Policy questions in the Solow model} (8).
\begin{flatlist}
  \item \colorbox{keybg}{\strut\ blue boxes\ } --- results to remember (mostly exactly as on the slides).
  \item \colorbox{intbg}{\strut\ green bars\ } --- the idea in plain words.
  \item \colorbox{pitbg}{\strut\ orange bars\ } --- easy-to-make mistakes.
  \item Dashed boxes marked \emph{Beyond the slides} --- worked checks and extra steps that are \emph{not} on the slides. Skip them if you only want the lecture content.
  \item In derivations, the grey italic text at the right edge says \emph{what was done} on that line.
\end{flatlist}
\end{recap}

{\small\tableofcontents}
\newpage

% =========================================================================
\section{Notation and the Growth-Rate Toolkit}
\slides{L1 2--3}

\subsection{Three ways to measure every variable}

The model is written three times over, once for each way of measuring a quantity:

\begin{center}
\renewcommand{\arraystretch}{1.35}
\begin{tabular}{@{}lllp{5.4cm}@{}}
\toprule
Written as & Example & Definition & Meaning \\
\midrule
Upper case & $Y_t, K_t$ & --- & aggregate (whole economy) \\
Lower case & $y_t, k_t$ & $y_t = Y_t/L_t$ & per capita (per worker) \\
Hat & $\yh_t, \kh_t$ & $\yh_t = Y_t/(A_tL_t)$ & per \emph{effective} unit of labour \\
\bottomrule
\end{tabular}
\end{center}

$A_tL_t$ is the number of \emph{effective} labour units: workers $L_t$ weighted by how productive each one is, $A_t$.

\begin{center}
\renewcommand{\arraystretch}{1.2}
\begin{tabular}{@{}ll@{\qquad}ll@{}}
\toprule
\multicolumn{4}{@{}l}{\emph{Variables and parameters (L1 slide 3)}}\\
\midrule
$Y_t$ & GDP (output) & $A_t$ & productivity (labour-augmenting technology) \\
$C_t$ & consumption & $x$ & growth rate of technology \\
$I_t$ & investment & $w_t$ & wage rate \\
$K_t$ & capital stock & $r_t$ & rental rate of capital \\
$L_t$ & population / labour & $n$ & population growth rate \\
$s$ & saving rate & $\delta$ & depreciation rate \\
\bottomrule
\end{tabular}
\end{center}
The parameters $s, n, x, \delta$ are \textbf{constant and exogenous}: the model takes them as given.

\subsection{Dots and growth rates}

A dot is a time derivative, $\dot X_t \equiv dX_t/dt$: how much $X$ changes per unit of time.
Dividing by $X$ gives the \textbf{growth rate}, the change as a \emph{fraction} of the current level:
\[
\gr{X_t} = \frac{dX_t/dt}{X_t}.
\]
(E.g.\ $\dot L_t/L_t = 0.01$ means population grows 1\% per year.)

\topic{Why a growth rate is the time derivative of a log}
\begin{align*}
\frac{d\ln X_t}{dt} &= \frac{d\ln X_t}{dX_t}\cdot\frac{dX_t}{dt} \why{chain rule}\\
&= \frac{1}{X_t}\cdot \dot X_t = \gr{X_t}. \why{$d\ln X/dX = 1/X$}
\end{align*}
This one fact turns every growth-rate calculation into ``take logs, then differentiate''.

\topic{The three growth-rate rules}
Take logs, then use $\ln(XY) = \ln X + \ln Y$, $\ln(X/Y) = \ln X - \ln Y$, $\ln X^a = a\ln X$, and differentiate with respect to $t$:
\begin{align*}
\frac{d\ln(XY)}{dt} &= \frac{d\ln X}{dt} + \frac{d\ln Y}{dt} \quad\Longrightarrow\quad \frac{(\dot{XY})}{XY} = \gr{X} + \gr{Y}, \why{product}\\
\frac{d\ln(X/Y)}{dt} &= \frac{d\ln X}{dt} - \frac{d\ln Y}{dt} \quad\Longrightarrow\quad \frac{(\dot{X/Y})}{X/Y} = \gr{X} - \gr{Y}, \why{ratio}\\
\frac{d\ln(X^a)}{dt} &= a\,\frac{d\ln X}{dt} \quad\Longrightarrow\quad \frac{(\dot{X^a})}{X^a} = a\,\gr{X}. \why{power}
\end{align*}

\begin{keyresult}[Growth-rate rules]
Products $\to$ \textbf{add} growth rates; ratios $\to$ \textbf{subtract} them; powers $\to$ \textbf{multiply} by the exponent.
\end{keyresult}

\subsection{Exogenous growth of population and technology}
Population and technology grow at constant rates:
\[
\dot L_t = nL_t \iff \gr{L_t} = n, \qquad\qquad \gr{A_t} = x.
\]
Solving the first one (the second is identical with $x$ in place of $n$):
\begin{align*}
\frac{d\ln L_t}{dt} &= n \why{growth rate = derivative of the log}\\
\ln L_t &= \ln L_0 + nt \why{integrate from $0$ to $t$}\\
L_t &= L_0e^{nt}. \why{exponentiate}
\end{align*}
Likewise $A_t = A_0e^{xt}$. By the product rule, effective labour grows at the \emph{sum} of the two rates:
\[
\frac{(\dot{A_tL_t})}{A_tL_t} = x + n \qquad\Longrightarrow\qquad A_tL_t = A_0L_0e^{(n+x)t}.
\]

% =========================================================================
\section{Production and Factor Prices}
\slides{L1 4--8}

\subsection{The production function}
Output is produced from effective labour and capital, $Y_t = F(A_tL_t, K_t)$, with three properties:
\begin{flatnum}
  \item \textbf{Constant returns to scale (CRS).} Double both inputs and output doubles:
  \[ F(A_t\lambda L_t, \lambda K_t) = \lambda F(A_tL_t, K_t) = \lambda Y_t \quad\text{for any } \lambda > 0. \]
  \item \textbf{Positive but diminishing marginal products.} More of an input raises output, but by less and less:
  \[ F_L > 0,\quad F_K > 0, \qquad F_{LL} < 0,\quad F_{KK} < 0. \]
  \item \textbf{Inada conditions.}
  \[ \lim_{L_t\to 0}F_L = \infty,\quad \lim_{L_t\to\infty}F_L = 0, \qquad \lim_{K_t\to 0}F_K = \infty,\quad \lim_{K_t\to\infty}F_K = 0. \]
  As an input becomes scarce its marginal product becomes arbitrarily large; as it becomes abundant its marginal product goes to zero.
\end{flatnum}

\topic{Intensive form: one variable is enough}
Under CRS we can choose $\lambda = 1/(A_tL_t)$:
\[
\yh_t = \frac{Y_t}{A_tL_t} = F\!\left(1, \frac{K_t}{A_tL_t}\right) \equiv f(\kh_t).
\]
So output per effective worker depends \emph{only} on capital per effective worker. This is why the whole model can later be reduced to one equation in $\kh_t$.

\subsection{Cobb--Douglas production}
The slides use
\[
Y_t = (A_tL_t)^{1-\alpha}K_t^{\alpha}, \qquad 0 < \alpha < 1,
\]
which satisfies all three properties. Dividing by $A_tL_t$ and by $L_t$ gives the other two versions:
\begin{align*}
\yh_t = \frac{(A_tL_t)^{1-\alpha}K_t^{\alpha}}{A_tL_t} &= (A_tL_t)^{-\alpha}K_t^{\alpha} = \left(\frac{K_t}{A_tL_t}\right)^{\!\alpha} = \kh_t^{\alpha}, \why{divide by $A_tL_t$}\\[2pt]
y_t = \frac{(A_tL_t)^{1-\alpha}K_t^{\alpha}}{L_t} &= A_t^{1-\alpha}L_t^{-\alpha}K_t^{\alpha} = A_t^{1-\alpha}\left(\frac{K_t}{L_t}\right)^{\!\alpha} = A_t^{1-\alpha}k_t^{\alpha}. \why{divide by $L_t$}
\end{align*}

\subsection{Factor prices and constant income shares}
Competitive firms pay each factor the value of its marginal product:
\begin{align*}
w_t &= \pd{Y_t}{L_t} = (1-\alpha)A_t^{1-\alpha}L_t^{-\alpha}K_t^{\alpha} = (1-\alpha)\frac{Y_t}{L_t}, \why{differentiate w.r.t.\ $L_t$}\\
r_t &= \pd{Y_t}{K_t} = \alpha(A_tL_t)^{1-\alpha}K_t^{\alpha-1} = \alpha\frac{Y_t}{K_t}. \why{differentiate w.r.t.\ $K_t$}
\end{align*}
In per-effective-worker terms (these become equations (5) and (6) below):
\[
r_t = \alpha\kh_t^{\alpha-1}, \qquad w_t = (1-\alpha)A_t\kh_t^{\alpha}.
\]
Multiplying $w_t$ by $L_t$ and $r_t$ by $K_t$ gives the income shares:
\[
\frac{w_tL_t}{Y_t} = 1-\alpha, \qquad \frac{r_tK_t}{Y_t} = \alpha, \qquad\text{so}\quad w_tL_t + r_tK_t = Y_t.
\]

\begin{keyresult}[Constant factor shares]
Labour's share of income is $1-\alpha$ and capital's share is $\alpha$, whatever the values of $K$, $L$, $A$. The data show roughly this: U.S.\ labour's share (labour compensation / total income, annual, 1929--2024) fluctuates around $62.8\%$ ($64.2\%$ post-1960). The standard calibration is
\[ 1-\alpha = 2/3 \quad\Longrightarrow\quad \alpha = 1/3. \]
\end{keyresult}

% =========================================================================
\section{The Equations of the Model}
\slides{L1 9--11}

\subsection{The model in three units}
\begin{center}
\renewcommand{\arraystretch}{1.3}
\begin{tabular}{@{}lllll@{}}
\toprule
 & Aggregate & Per capita & Per effective worker & \\
\midrule
Production & $Y = (AL)^{1-\alpha}K^{\alpha}$ & $y = A^{1-\alpha}k^{\alpha}$ & $\yh = \kh^{\alpha}$ & (1) \\
Resource constraint & $Y = C + I$ & $y = c + i$ & $\yh = \ch + \ih$ & (2) \\
Saving = investment & $I = sY$ & $i = sy$ & $\ih = s\yh$ & (3) \\
Capital accumulation & $\dot K = I - \delta K$ & $\dot k = i - (\delta+n)k$ & $\dk = \ih - \gam\kh$ & (4) \\
\bottomrule
\end{tabular}
\end{center}
plus the factor prices $r = \alpha\kh^{\alpha-1}$ \;(5) and $w = (1-\alpha)A\kh^{\alpha}$ \;(6). Time subscripts are dropped.

Rows (1)--(3) are just the aggregate equation divided by $L$ or by $AL$. Row (4) needs more care, because \emph{the time derivative of a ratio is not the ratio of time derivatives}.

\subsection{Deriving the capital-accumulation equations}

\begin{pitfall}
$\dfrac{\dot K}{L} = \dfrac{I}{L} - \delta\dfrac{K}{L} = i - \delta k$, but this is \textbf{not} $\dot k$. Capital per worker also falls when population grows, because the same $K$ is shared among more people. Always go through growth rates.
\end{pitfall}

\topic{Per capita: $k = K/L$}
\begin{align*}
\gr{k} &= \gr{K} - \gr{L} \why{ratio rule}\\
&= \frac{I - \delta K}{K} - n \why{use $\dot K = I - \delta K$ and $\dot L/L = n$}\\
\dot k &= \frac{I - \delta K}{K}\cdot\frac{K}{L} - nk \why{multiply both sides by $k = K/L$}\\
&= \frac{I}{L} - \delta k - nk = i - (\delta+n)k. \why{simplify}
\end{align*}

\topic{Per effective worker: $\kh = K/(AL)$}
\begin{align*}
\gr{\kh} &= \gr{K} - \gr{A} - \gr{L} \why{ratio + product rules}\\
&= \frac{I - \delta K}{K} - x - n \why{use $\dot K$, $\dot A/A = x$, $\dot L/L = n$}\\
\dk &= \frac{I - \delta K}{K}\cdot\frac{K}{AL} - (x+n)\kh \why{multiply by $\kh = K/(AL)$}\\
&= \frac{I}{AL} - \delta\kh - (x+n)\kh = \ih - \gam\kh. \why{simplify}
\end{align*}

\begin{intuition}
$(\delta+n+x)\kh$ is \textbf{break-even investment}: the investment needed just to keep $\kh$ constant. It must replace worn-out capital ($\delta$), equip the new workers ($n$), and keep up with each worker becoming more productive ($x$).
\end{intuition}

\subsection{The key dynamic equation}
From (2) and (3), consumption is what is not saved:
\[
\ch_t = (1-s)\yh_t. \tag{7}
\]
Substituting (1) and (3) into (4):
\begin{keyresult}[The Solow equation]
\[
\dk_t = s\yh_t - \gam\kh_t = s\kh_t^{\alpha} - \gam\kh_t \tag{8}
\]
Actual investment minus break-even investment. It is a \textbf{nonlinear first-order differential equation} in $\kh_t$. The goal (Section~\ref{sec:ode}) is a solution relating $\kh_t$ to time, the parameters $(s,\delta,n,x,\alpha)$, and the initial value $\kh_0$.
\end{keyresult}

% =========================================================================
\section{Dynamics and the Steady State}
\slides{L1 12--13; EV 3, 14}

\subsection{The Solow diagram}
Plot the two terms of (8) against $\kh$. Because $0 < \alpha < 1$, actual investment $s\kh^{\alpha}$ is concave; break-even investment $\gam\kh$ is a straight line through the origin. $\dk$ is the vertical gap between them.

\begin{figure}[ht]
\centering
\begin{tikzpicture}[xscale=0.55,yscale=5.2]
  \draw[ecaxis] (0,0) -- (12.5,0);
  \draw[ecaxis] (0,0) -- (0,0.95);
  \node[ecnode,anchor=west] at (12.6,0) {$\kh$};
  \draw[budget] plot[domain=0.001:12,samples=80] (\x,{0.3*\x^(1/3)});
  \draw[budget,dashed] (0,0) -- (11,0.88);
  \node[ecnode,anchor=west] at (12,{0.3*12^(1/3)}) {$\ih = s\kh^{\alpha}$};
  \node[ecnode,anchor=west] at (11,0.88) {$\gam\kh$};
  \draw[densely dotted] (7.26,0) -- (7.26,0.581) node[bundle] {};
  \node[ecnode,below] at (7.26,0) {$\kh^{*}$};
  \draw[-{Stealth[length=4pt]},line width=0.8pt] (2.0,0.03) -- (4.5,0.03);
  \draw[-{Stealth[length=4pt]},line width=0.8pt] (11.5,0.03) -- (9.3,0.03);
  \node[ecnode,align=center] at (3.2,0.63) {$s\kh^{\alpha} > \gam\kh$\\$\implies \dk > 0$};
  \node[ecnode,align=center] at (10.3,0.35) {$s\kh^{\alpha} < \gam\kh$\\$\implies \dk < 0$};
\end{tikzpicture}
\caption{The Solow diagram ($s=0.3$, $\alpha=1/3$, $\delta+n+x=0.08$). Arrows show the direction in which $\kh$ moves.}
\end{figure}

\begin{center}
\renewcommand{\arraystretch}{1.3}
\begin{tabular}{@{}lll@{}}
\toprule
Where $\kh$ is & Comparison & Consequence \\
\midrule
$\kh_t < \kh^*$ & $\ih_t > \gam\kh_t$ & $\dk_t > 0$: $\kh$ rises \\
$\kh_t = \kh^*$ & $\ih_t = \gam\kh_t$ & $\dk_t = 0$: steady state \\
$\kh_t > \kh^*$ & $\ih_t < \gam\kh_t$ & $\dk_t < 0$: $\kh$ falls \\
\bottomrule
\end{tabular}
\end{center}

\subsection{Solving for the steady state}
Set $\dk = 0$ in (8) and solve for $\kh > 0$:
\begin{align*}
0 &= s(\kh^*)^{\alpha} - \gam\kh^* \why{steady state: $\dk = 0$}\\
(\kh^*)^{\alpha-1} &= \frac{\delta+n+x}{s} \why{divide by $s\kh^*$}\\
\kh^* &= \left(\frac{s}{\delta+n+x}\right)^{\frac{1}{1-\alpha}}. \why{raise to the power $\tfrac{1}{\alpha-1}$}
\end{align*}
The other steady-state values follow:
\[
\yh^* = (\kh^*)^{\alpha} = \left(\frac{s}{\delta+n+x}\right)^{\frac{\alpha}{1-\alpha}}, \qquad
\ch^* = (1-s)\yh^*, \qquad
\left(\frac{Y}{L}\right)^{\!*} = A\yh^* = A\left(\frac{s}{\delta+n+x}\right)^{\frac{\alpha}{1-\alpha}}.
\]

\begin{keyresult}[Steady state]
\[
\kh^* = \left(\frac{s}{\delta+n+x}\right)^{\frac{1}{1-\alpha}}, \qquad
\left(\frac{Y}{L}\right)^{\!*} = A\left(\frac{s}{\delta+n+x}\right)^{\frac{\alpha}{1-\alpha}}.
\]
Long-run income depends on technology $A$ and on the parameters that determine steady-state capital intensity.
\end{keyresult}

\subsection{Convergence to the steady state}
Divide (8) by $\kh$ to get the growth rate of $\kh$:
\[
\gr{\kh} = s\kh^{\alpha-1} - \gam, \qquad
\frac{\partial(\dk/\kh)}{\partial\kh} = (\alpha-1)s\kh^{\alpha-2} < 0 \quad(\text{because } \alpha < 1).
\]
The growth rate falls steadily as $\kh$ rises: it is very large when $\kh$ is close to $0$, it approaches $-\gam < 0$ when $\kh$ is very large, and it equals zero only at $\kh^*$. So it is positive below $\kh^*$ and negative above it.

\begin{keyresult}[Global stability]
The system converges to the steady state from \emph{any} starting value $\kh_0 > 0$. $\kh$ grows only during the transition; once $\kh^*$ is reached, growth in capital per effective worker stops.
\end{keyresult}

% =========================================================================
\section{Balanced Growth: The Long-Run Facts}
\slides{EV 2--6; L1 14}

\emph{Question (EV slide 2): does the model reproduce the balanced-growth properties of developed economies?}

\subsection{Output per person grows at the rate of technological progress}
Output per person is $y = A\yh = A\kh^{\alpha}$. By the growth-rate rules:
\[
\gr{y} = \gr{A} + \gr{\yh} = x + \alpha\gr{\kh}.
\]
As $\kh \to \kh^*$, $\dk/\kh \to 0$. \textbf{Hence $\dot y/y \to x$.}

\subsection{Consumption and investment grow with output}
The resource constraint and the constant saving rate give $C = (1-s)Y$ and $I = sY$. Since $s$ is constant,
\[
\gr{C} = \gr{Y}, \qquad \gr{I} = \gr{Y}.
\]
In the long run output, consumption and investment \emph{per person} all grow at rate $x$. This broadly describes developed economies; in many developing economies, however, saving rates rise during development.

\subsection{Capital grows faster than labour}
From $\kh = K/(AL)$:
\[
\gr{K} = \gr{A} + \gr{L} + \gr{\kh} = x + n + \gr{\kh}.
\]
On the balanced growth path $\dk/\kh = 0$, so
\[
\gr{K} = n + x \qquad\text{and}\qquad \gr{K} - \gr{L} = x.
\]

\subsection{The return to capital is constant}
The net return to capital is
\[
r - \delta = \alpha(AL)^{1-\alpha}K^{\alpha-1} - \delta = \alpha\kh^{\alpha-1} - \delta.
\]
On the balanced growth path $\kh$ is constant, so $r - \delta$ is constant. $K$ and $AL$ both grow, but their effects on the marginal product of capital offset each other.

\subsection{Summary of long-run growth rates}
\begin{center}
\renewcommand{\arraystretch}{1.25}
\begin{tabular}{@{}llc@{}}
\toprule
Variable & & Long-run growth rate \\
\midrule
Per effective worker & $\kh, \yh, \ch, \ih$ & $0$ \\
Per capita & $y, c, i, k$, and the wage $w$ & $x$ \\
Aggregates & $Y, C, I, K$ & $n + x$ \\
Rental rate & $r$ (and $r - \delta$) & $0$ \\
Ratios & $K/Y,\; C/Y,\; I/Y$ & $0$ \\
\bottomrule
\end{tabular}
\end{center}

\begin{recap}
In the long run everything ``per effective worker'' is constant, everything ``per person'' grows at the technology rate $x$, and aggregates grow at $n + x$. That is the model's \emph{balanced growth path}.
\end{recap}

% =========================================================================
\section{Transitional Growth and Convergence}
\slides{EV 7--13}

\subsection{Long-run growth vs.\ transitional growth}
Combining $\dot y/y = x + \alpha\,\dk/\kh$ with the growth rate of $\kh$:
\[
\gr{y} = \underbrace{x}_{\text{long-run growth}} + \underbrace{\alpha\gr{\kh}}_{\text{transitional growth}},
\qquad \gr{\kh} = s\kh^{\alpha-1} - \delta - n - x, \qquad \frac{\partial(\dk/\kh)}{\partial\kh} < 0.
\]
\begin{intuition}[Key mechanism]
An economy with little capital per effective worker has a \emph{high marginal product of capital}, so it grows fast during its transition to the steady state. As capital accumulates, diminishing returns set in and the extra growth fades.
\end{intuition}

\subsection{Conditional convergence, no absolute convergence, divergence}
\begin{flatlist}
  \item \textbf{Conditional convergence.} Suppose a country is poor because its initial $K/L$ is low, while $s, n, x, \delta$ and the technology $A$ are the same as elsewhere. Then
  \[ \kh \text{ low} \implies Y/L \text{ low} \quad\text{and}\quad \dk/\kh \text{ high} \implies \dot y/y \text{ high}. \]
  Countries converge to the same steady state when the parameters that determine it are the same.

  \item \textbf{Why absolute convergence need not occur.} If a country is poor because its technology $A$ is low, its $\kh = K/(AL)$ may be \emph{high}. Then
  \[ Y/L \text{ low}, \qquad \kh \text{ high}, \qquad \dk/\kh \text{ low}, \qquad \dot y/y \text{ low}. \]
  Differences in $s$, $n$ or $x$ also produce different steady states. A poor country need not grow faster than a rich one.

  \item \textbf{Divergence.} Take two countries that start with the same income per worker. If one saves more, it converges to a higher steady-state \emph{level}; if one has faster technological progress, its long-run \emph{growth rate} is higher. Their incomes diverge.
\end{flatlist}

\begin{keyresult}[Prediction]
The Solow model predicts \textbf{conditional} convergence, not convergence among all countries regardless of their characteristics. Empirically: convergence within groups of similar countries or regions, persistent or widening gaps across dissimilar ones.
\end{keyresult}

\subsection{A numerical example}
Let $s = 0.25$, $\alpha = 1/3$, $\delta = 0.05$, $n = 0.01$, $x = 0.02$ (so $\delta+n+x = 0.08$), and start at $\kh_0 = 1.8$. Initial growth of output per worker:
\begin{align*}
\gr{y} &= x + \alpha\left[s\kh_0^{\alpha-1} - \gam\right] \\
&= 0.02 + \tfrac{1}{3}\left[0.25\times 1.8^{-2/3} - 0.08\right] \why{plug in}\\
&= 0.02 + \tfrac{1}{3}\left[0.169 - 0.080\right] \approx 0.050. \why{$1.8^{-2/3} \approx 0.676$}
\end{align*}
Growth starts near \textbf{5\%}, but the transitional part fades as capital accumulates:

\begin{center}
\small
\begin{tabular}{@{}rccc@{\qquad}rccc@{}}
\toprule
$t$ & $\kh_t$ & $\dk_t/\kh_t$ & $\dot y_t/y_t$ & $t$ & $\kh_t$ & $\dk_t/\kh_t$ & $\dot y_t/y_t$ \\
\midrule
0 & 1.8000 & -- & -- & 7 & 2.8267 & 0.0491 & 0.0364 \\
1 & 1.9601 & 0.0890 & 0.0497 & 8 & 2.9541 & 0.0451 & 0.0350 \\
2 & 2.1162 & 0.0796 & 0.0465 & 9 & 3.0764 & 0.0414 & 0.0338 \\
3 & 2.2678 & 0.0717 & 0.0439 & 10 & 3.1939 & 0.0382 & 0.0327 \\
4 & 2.4149 & 0.0648 & 0.0416 & 11 & 3.3066 & 0.0353 & 0.0318 \\
5 & 2.5571 & 0.0589 & 0.0396 & 12 & 3.4145 & 0.0326 & 0.0309 \\
6 & 2.6944 & 0.0537 & 0.0379 & 13 & 3.5178 & 0.0303 & 0.0301 \\
 & & & & 14 & 3.6166 & 0.0281 & 0.0294 \\
\bottomrule
\end{tabular}
\end{center}
Output growth falls quickly toward the long-run rate of 2\%.

\begin{extra}
The table can be reproduced with annual steps $\kh_{t+1} = \kh_t + s\kh_t^{\alpha} - 0.08\,\kh_t$; e.g.\ $\kh_1 = 1.8 + 0.25(1.8)^{1/3} - 0.144 = 1.9601$. The columns are $\dk_t/\kh_t = (\kh_t - \kh_{t-1})/\kh_{t-1}$ and $\dot y_t/y_t = x + \alpha\,\dk_t/\kh_t$.
\end{extra}

\subsection{Limits of transitional growth}
Because $\dk/\kh \to 0$, the model has difficulty sustaining a growth miracle for several decades:
\begin{flatlist}
  \item sustained growth near 5\% eventually requires technological progress near 5\%;
  \item a prolonged growth disaster similarly requires technological progress near zero or below zero.
\end{flatlist}
\begin{keyresult}
Capital accumulation can generate large \emph{short-run} growth differences, but technology determines \emph{long-run} growth.
\end{keyresult}

% =========================================================================
\section{Explaining International Income Differences}
\slides{EV 14--17}

\emph{Question (EV slide 2): can plausible differences in saving, population growth and technology explain international income gaps?}

Start from steady-state income per worker (Section 4.2):
\[
y^* = A\left(\frac{s}{\delta+n+x}\right)^{\frac{\alpha}{1-\alpha}}, \qquad \frac{\alpha}{1-\alpha} = \frac{1/3}{2/3} = \frac12.
\]
Can we get a \textbf{tenfold} gap, $y_R/y_P = 10$ (R = rich, P = poor)? Change one thing at a time.

\subsection{Saving rates}
Same $A, n, x, \delta, \alpha$; different $s$. Then
\begin{align*}
\frac{y_R}{y_P} &= \left(\frac{s_R}{s_P}\right)^{\frac{\alpha}{1-\alpha}} = 10 \why{only $s$ differs}\\
\frac{s_R}{s_P} &= 10^{(1-\alpha)/\alpha} = 10^2 = 100. \why{raise to the power $\tfrac{1-\alpha}{\alpha} = 2$}
\end{align*}
If $s_R = 0.25$, then $s_P = 0.25/100 = 0.0025$.

\textbf{Assessment.} Poor countries do tend to save less, but a 0.25\% saving rate is implausibly low. Saving differences explain only part of the gap.

\subsection{Population growth}
Same $A$ and $s$; different $n$. Then
\begin{align*}
\frac{y_R}{y_P} &= \left(\frac{\delta+n_P+x}{\delta+n_R+x}\right)^{\frac{\alpha}{1-\alpha}} = 10 \why{only $n$ differs}\\
\delta + n_P + x &= 100(\delta + n_R + x) \why{square both sides}\\
n_P &= 100(\delta + n_R + x) - \delta - x = 100(0.08) - 0.07 = 7.93. \why{$\delta=0.05$, $n_R=0.01$, $x=0.02$}
\end{align*}
That is about \textbf{793\% per year}.

\textbf{Assessment.} Population-growth differences of plausible size cannot explain a tenfold income gap.

\subsection{Technology}
All other parameters the same. Then
\[
\frac{y_R}{y_P} = \frac{A_R}{A_P} = 10.
\]
\textbf{Assessment.} Differences in technology / total factor productivity can, in accounting terms, generate large income differences. \emph{Limitation:} the basic Solow model treats the level of technology as exogenous, so it does not explain why technology differs across countries.

\begin{center}
\renewcommand{\arraystretch}{1.3}
\begin{tabular}{@{}lll@{}}
\toprule
Source of the gap & Required for $y_R/y_P = 10$ & Plausible? \\
\midrule
Saving rate $s$ & $s_R/s_P = 100$ (e.g.\ $s_P = 0.25\%$) & No \\
Population growth $n$ & $n_P \approx 793\%$ per year & No \\
Technology $A$ & $A_R/A_P = 10$ & Yes, but $A$ is unexplained \\
\bottomrule
\end{tabular}
\end{center}

\begin{intuition}
Because of diminishing returns to capital, steady-state income responds only weakly to saving and population growth (elasticity $\alpha/(1-\alpha) = 1/2$). Huge differences in these parameters are needed to produce modest income differences.
\end{intuition}

% =========================================================================
\section{Solving the Solow Differential Equation}\label{sec:ode}
\slides{DE 1--14}

\begin{recap}[Roadmap of this section]
(1) Learn to solve a \emph{linear} first-order ODE. \;
(2) Turn the (nonlinear) Solow equation into a linear one and solve it \emph{exactly}. \;
(3) Solve it \emph{approximately} by log-linearizing. \;
(4) Translate capital convergence into \emph{income} convergence. \;
(5) Turn that into a cross-country \emph{growth regression}. \;
(6) Compute the \emph{speed} of convergence and the half-life.
\end{recap}

\subsection{Linear first-order differential equations}
A first-order differential equation relates a variable, its time derivative, and possibly time: $\dot z_t = g(z_t, t)$. A \textbf{linear} one has the form
\[
\dot z_t + a(t)z_t = b(t).
\]
Together with an initial condition $z_0$, it determines the path $z_t$. Take constant coefficients $a > 0$ and $b$:
\[
\dot z_t + az_t = b.
\]

\topic{Solving it with an integrating factor}
\step{1}{Multiply both sides by $e^{at}$:}
\[ e^{at}\dot z_t + ae^{at}z_t = be^{at}. \]
\step{2}{The left-hand side is the derivative of a product (product rule: $\tfrac{d}{dt}(e^{at}z_t) = e^{at}\dot z_t + ae^{at}z_t$):}
\[ \frac{d}{dt}\left(e^{at}z_t\right) = be^{at}. \]
\step{3}{Integrate both sides over time ($C$ is an arbitrary constant):}
\[ \int\frac{d}{dt}\left(e^{at}z_t\right)dt = \int be^{at}\,dt \quad\Longrightarrow\quad e^{at}z_t = \frac{b}{a}e^{at} + C. \]
\step{4}{Divide by $e^{at}$ --- the family of all solutions:}
\[ z_t = \frac{b}{a} + Ce^{-at}. \]
\step{5}{Pin down $C$ with the initial condition. At $t = 0$, $e^{-a\cdot 0} = 1$, so}
\[ z_0 = \frac{b}{a} + C \quad\Longrightarrow\quad C = z_0 - \frac{b}{a}. \]

\begin{keyresult}[Solution of $\dot z_t + az_t = b$]
\[
z_t = \frac{b}{a} + \left(z_0 - \frac{b}{a}\right)e^{-at} \qquad\text{or, with } z^* \equiv b/a, \qquad z_t = z^* + (z_0 - z^*)e^{-at}.
\]
\end{keyresult}

\topic{What the solution says}
\begin{flatlist}
  \item The gap $z_t - z^*$ shrinks at exponential rate $a$. Differentiating with respect to time:
  \[ \frac{d}{dt}(z_t - z^*) = -a(z_0 - z^*)e^{-at} = -a(z_t - z^*). \]
  \item If $z_0 < z^*$ the path rises toward $z^*$; if $z_0 > z^*$ it falls.
  \item At time $t$, the fraction of the initial gap that has been closed is $1 - e^{-at}$.
\end{flatlist}

\subsection{The exact Solow transition path}
With Cobb--Douglas, the Solow equation
\[
\dk_t = s\kh_t^{\alpha} - \gam\kh_t, \qquad 0 < \alpha < 1,
\]
is \emph{nonlinear} because of the $\kh_t^{\alpha}$ term, so the formula above does not apply directly. Its positive steady state is $\kh^* = [s/\gam]^{1/(1-\alpha)}$. It is a \textbf{Bernoulli equation}: a change of variable makes it linear.

\step{1}{For $\kh_t > 0$, divide the equation by $\kh_t^{\alpha}$ (after moving the break-even term to the left):}
\[ \kh_t^{-\alpha}\dk_t + \gam\kh_t^{1-\alpha} = s. \]
\step{2}{Define the new variable $z_t \equiv \kh_t^{1-\alpha}$. By the chain rule,}
\[ \dot z_t = (1-\alpha)\kh_t^{-\alpha}\dk_t \quad\Longrightarrow\quad \kh_t^{-\alpha}\dk_t = \frac{\dot z_t}{1-\alpha}. \]
\step{3}{Substitute into Step 1 and multiply by $(1-\alpha)$:}
\[ \dot z_t + \underbrace{(1-\alpha)\gam}_{a}\,z_t = \underbrace{(1-\alpha)s}_{b}. \]
This is \textbf{linear in $z_t$}, exactly the form we just solved.

\step{4}{Its steady state is}
\[ z^* = \frac{b}{a} = \frac{(1-\alpha)s}{(1-\alpha)\gam} = \frac{s}{\delta+n+x} = (\kh^*)^{1-\alpha}. \]
\step{5}{Apply the linear solution, then undo the substitution ($\kh_t = z_t^{1/(1-\alpha)}$, $z_0 = \kh_0^{1-\alpha}$).}

\begin{keyresult}[Exact Solow path]
\[
z_t = z^* + (z_0 - z^*)e^{-\beta t}, \qquad \beta \equiv (1-\alpha)(\delta+n+x),
\]
\[
\kh_t = \left[z^* + \left(\kh_0^{1-\alpha} - z^*\right)e^{-\beta t}\right]^{\frac{1}{1-\alpha}}.
\]
\end{keyresult}

\topic{Exact and approximate convergence}
The exact solution says that the gap in $z_t = \kh_t^{1-\alpha}$ closes at the rate $\beta = (1-\alpha)(\delta+n+x)$. The path for $\kh_t$ itself is nonlinear. Near the steady state, however, its \emph{logarithm} approximately follows the same exponential convergence law (next subsection).

\emph{Why bother with an approximation?} Log output per effective worker is proportional to log capital per effective worker ($\ln\yh = \alpha\ln\kh$), which makes the log version useful for measuring \emph{income} convergence.

\begin{extra}[Beyond the slides --- numerical check]
With the Section 6.3 numbers ($s = 0.25$, $\alpha = 1/3$, $\delta+n+x = 0.08$, $\kh_0 = 1.8$):
$\beta = \tfrac23(0.08) \approx 0.0533$, $z^* = 0.25/0.08 = 3.125$, $z_0 = 1.8^{2/3} \approx 1.4798$. At $t = 14$:
\[
z_{14} = 3.125 - 1.6452\,e^{-0.0533\times 14} \approx 3.125 - 1.6452\times 0.4739 \approx 2.3453, \qquad \kh_{14} = 2.3453^{3/2} \approx 3.59,
\]
close to the table's $3.617$ (the table uses annual steps, not continuous time).
\end{extra}

\subsection{Log-linearizing around the steady state}
\step{1}{Work in logs. Let $\tilde k_t \equiv \ln\kh_t$, so $\kh_t = e^{\tilde k_t}$. Then}
\begin{align*}
\dot{\tilde k}_t = \frac{\dk_t}{\kh_t} &= s\kh_t^{\alpha-1} - \gam \why{growth rate of $\kh$}\\
&= se^{-(1-\alpha)\tilde k_t} - \gam \;\equiv\; h(\tilde k_t). \why{$\kh^{\alpha-1} = e^{(\alpha-1)\tilde k}$}
\end{align*}
\step{2}{First-order Taylor expansion of $h$ around $\tilde k^* = \ln\kh^*$:}
\[
\dot{\tilde k}_t = h(\tilde k_t) \simeq \underbrace{h(\tilde k^*)}_{=\,\dot{\tilde k}^* = 0} + h'(\tilde k^*)(\tilde k_t - \tilde k^*).
\]
The first term is zero because $\tilde k$ does not change in the steady state. So we only need the slope $h'(\tilde k^*)$.

\step{3}{Differentiate $h$ and evaluate at the steady state:}
\begin{align*}
h'(\tilde k) &= -(1-\alpha)se^{-(1-\alpha)\tilde k} \why{chain rule}\\
se^{-(1-\alpha)\tilde k^*} &= \delta + n + x \why{steady state: $s(\kh^*)^{\alpha-1} = \gam$}\\
h'(\tilde k^*) &= -(1-\alpha)(\delta+n+x) \equiv -\beta. \why{combine}
\end{align*}
\step{4}{Substitute into the Taylor approximation:}
\[
\dot{\tilde k}_t \simeq -\beta(\tilde k_t - \tilde k^*), \qquad \beta = (1-\alpha)(\delta+n+x).
\]

\subsection{Solving the log-linear equation}
Rearrange it into the familiar linear form:
\[
\dot{\tilde k}_t + \beta\tilde k_t = \beta\tilde k^*.
\]
Match it with $\dot z_t + az_t = b$:
\[
z_t \longleftrightarrow \tilde k_t, \qquad a \longleftrightarrow \beta, \qquad b \longleftrightarrow \beta\tilde k^*, \qquad\text{so}\quad \frac{b}{a} = \tilde k^*.
\]
Applying $z_t = b/a + (z_0 - b/a)e^{-at}$:
\begin{keyresult}[Log-linear path]
\[
\tilde k_t \simeq \tilde k^* + (\tilde k_0 - \tilde k^*)e^{-\beta t}
\qquad\Longleftrightarrow\qquad
\tilde k_t - \tilde k_0 \simeq \left(1 - e^{-\beta t}\right)(\tilde k^* - \tilde k_0).
\]
Log capital closes its gap to the steady state at (approximately) rate $\beta$ per year.
\end{keyresult}

\begin{figure}[ht]
\centering
\begin{tikzpicture}[xscale=0.13,yscale=0.8,
  declare function={bet=0.0533333; ze(\t)=3.125-1.645273*exp(-bet*\t);
                    kl(\t)=exp(ln(5.524272)+(ln(1.8)-ln(5.524272))*exp(-bet*\t));}]
  \draw[ecaxis] (0,0) -- (84,0);
  \draw[ecaxis] (0,0) -- (0,6.4);
  \node[ecnode,anchor=west] at (84,0) {$t$};
  \node[ecnode,anchor=south] at (0,6.4) {$\kh_t$};
  \draw[densely dotted] (0,5.524) -- (80,5.524);
  \node[ecnode,left] at (0,5.524) {$\kh^{*}=5.52$};
  \node[ecnode,left] at (0,1.8) {$\kh_0=1.8$};
  \draw[budget] plot[domain=0:80,samples=80] (\x,{ze(\x)^1.5});
  \draw[budget,dashed] plot[domain=0:80,samples=80] (\x,{kl(\x)});
  \foreach \t in {20,40,60} \node[ecnode,below] at (\t,0) {$\t$};
  \node[ecnode,anchor=east] (ex) at (14,4.7) {exact};
  \draw[thin,black!55] (ex.east) -- (18,{ze(18)^1.5});
  \node[ecnode,anchor=west] at (38,3.2) {log-linear approx.};
  \draw[thin,black!55] (38,3.2) -- (34,{kl(34)});
\end{tikzpicture}
\caption{Exact path vs.\ log-linear approximation ($\kh_0 = 1.8$, $\kh^* = 5.52$, $\beta = 0.0533$). Far from the steady state the approximation understates the early catch-up; the two paths merge as $\kh_t \to \kh^*$. \emph{(Illustration, not on the slides.)}}
\end{figure}

\subsection{From capital convergence to income convergence}
Let $y_t \equiv Y_t/L_t = A_t\yh_t$ with $\yh_t = \kh_t^{\alpha}$.
\begin{align*}
\ln y_t &= \ln A_t + \alpha\tilde k_t, \qquad \ln\yh^* = \alpha\tilde k^* \why{take logs; $\tilde k = \ln\kh$}\\
\ln y_t - \ln y_0 &= (\ln A_t - \ln A_0) + \alpha(\tilde k_t - \tilde k_0) \why{subtract the $t=0$ version}\\
&= xt + \alpha(\tilde k_t - \tilde k_0) \why{$A$ grows at rate $x$}\\
&\simeq xt + \left(1 - e^{-\beta t}\right)\alpha(\tilde k^* - \tilde k_0). \why{log-linear path}
\end{align*}
Finally, $\alpha\tilde k^* = \ln\yh^*$ and $\alpha\tilde k_0 = \ln\yh_0 = \ln y_0 - \ln A_0$, so
\begin{keyresult}[Income convergence]
\[
\ln y_t - \ln y_0 \simeq xt + \left(1 - e^{-\beta t}\right)\left[\ln\yh^* - (\ln y_0 - \ln A_0)\right].
\]
Holding fixed the determinants of $\yh^*$ (and $A_0$), a \textbf{lower initial income $y_0$ predicts faster transitional growth}. The approximation applies near the steady state.
\end{keyresult}

\subsection{From the Solow transition to a growth regression}
Let $i$ index countries and $T$ be the length of the observation interval.

\step{1}{Write the income-convergence equation for country $i$ at $t = T$ and divide by $T$:}
\[
\frac{\ln y_{i,T} - \ln y_{i,0}}{T} \simeq x + \frac{1 - e^{-\beta T}}{T}\left[\ln\yh_i^* - \ln y_{i,0} + \ln A_{i,0}\right].
\]
\step{2}{Define the convergence coefficient $\lambda_T \equiv \dfrac{1 - e^{-\beta T}}{T}$.}

\step{3}{Substitute the Cobb--Douglas steady state:}
\[
\ln\yh_i^* = \frac{\alpha}{1-\alpha}\ln s_i - \frac{\alpha}{1-\alpha}\ln(n_i + x + \delta).
\]
\step{4}{Collect terms:}
\begin{keyresult}[Growth regression]
\[
\frac{\ln y_{i,T} - \ln y_{i,0}}{T} \simeq c_i - \lambda_T\ln y_{i,0} + \lambda_T\frac{\alpha}{1-\alpha}\ln s_i - \lambda_T\frac{\alpha}{1-\alpha}\ln(n_i + x + \delta),
\qquad c_i \equiv x + \lambda_T\ln A_{i,0}.
\]
\end{keyresult}
For estimation, initial productivity may be related to observed country characteristics $Z_i$ (e.g.\ rule of law, life expectancy, inflation):
\[
c_i = c_0 + \gamma'Z_i + u_i.
\]
\textbf{Conditional convergence predicts a negative coefficient on initial income} ($-\lambda_T < 0$), holding the determinants of each country's steady state fixed.

\subsection{Speed of convergence and half-life}
In the basic model $\beta = (1-\alpha)(\delta+n+x)$. With $\alpha = 1/3$, $\delta = 0.05$, $n = 0.01$, $x = 0.02$ per year:
\[
\beta = \tfrac{2}{3}(0.08) \simeq 0.0533 \text{ per year}.
\]
The \textbf{half-life} $T_{1/2}$ is the time needed to close half of the initial log gap:
\begin{align*}
e^{-\beta T_{1/2}} &= \tfrac12 \why{half the gap remains}\\
T_{1/2} &= -\frac{\ln(1/2)}{\beta} = \frac{\ln 2}{\beta}. \why{take logs}
\end{align*}
\begin{keyresult}[Half-life]
\begin{align*}
\text{model:}\quad &\beta = 0.0533 &&\Longrightarrow\quad T_{1/2} = 0.693/0.0533 \approx 13 \text{ years},\\
\text{illustrative empirical rate:}\quad &\beta = 0.02 &&\Longrightarrow\quad T_{1/2} = 0.693/0.02 \approx 35 \text{ years}.
\end{align*}
The model's convergence is too fast compared with the data. The slower empirical rate motivates richer models of convergence.
\end{keyresult}

% =========================================================================
\section{Policy Analysis}
\slides{PO 1--8}

Three questions (PO slide 1): What saving rate maximizes steady-state consumption? How does more saving affect consumption along the transition? Which forms of foreign aid change income temporarily or permanently? Throughout we use
\[
\yh = \kh^{\alpha}, \qquad \ch = (1-s)\yh, \qquad \dk = s\kh^{\alpha} - \gam\kh, \qquad 0 < \alpha < 1.
\]

\subsection{The Golden Rule}
In a steady state, investment just covers break-even investment, $\ih^* = \gam\kh^*$, so
\begin{align*}
\ch^* &= \yh^* - \ih^* = (\kh^*)^{\alpha} - \gam\kh^*. \why{consumption = output $-$ investment}\\
\frac{d\ch^*}{d\kh^*} &= \alpha(\kh^*)^{\alpha-1} - \gam = 0 \why{maximize over $\kh^*$}\\
\kh_{GR} &= \left(\frac{\alpha}{\delta+n+x}\right)^{\frac{1}{1-\alpha}}. \why{solve}
\end{align*}

\topic{Which saving rate implements the Golden Rule?}
Put the two formulas side by side:
\[
\kh^*(s) = \left(\frac{s}{\delta+n+x}\right)^{\frac{1}{1-\alpha}}
\qquad\text{vs.}\qquad
\kh_{GR} = \left(\frac{\alpha}{\delta+n+x}\right)^{\frac{1}{1-\alpha}}.
\]
They have the same form, with $s$ in place of $\alpha$.
\begin{keyresult}[Golden Rule saving rate]
\[ s_{GR} = \alpha. \]
Under Cobb--Douglas, the saving rate that maximizes steady-state consumption equals capital's income share.
\end{keyresult}

\begin{figure}[ht]
\centering
\begin{tikzpicture}[xscale=0.42,yscale=1.75]
  \draw[ecaxis] (0,0) -- (15,0);
  \draw[ecaxis] (0,0) -- (0,2.75);
  \node[ecnode,anchor=west] at (15.1,0) {$\kh$};
  \draw[budget] plot[domain=0.001:14.3,samples=90] (\x,{\x^(1/3)});
  \node[ecnode,anchor=west] at (14.3,{14.3^(1/3)}) {$\yh = \kh^{\alpha}$};
  \draw[budget,dashed] (0,0) -- (14.3,1.144);
  \node[ecnode,anchor=west] at (14.3,1.144) {$\gam\kh$};
  \draw[budget,black!55] plot[domain=0.001:14.3,samples=90] (\x,{\x^(1/3)/3});
  \node[ecnode,anchor=west,black!70] at (14.3,{14.3^(1/3)/3}) {$\alpha\kh^{\alpha}$};
  \draw[thin] (2.0,1.522) -- (13.5,2.442);
  \draw[densely dotted] (8.505,0) -- (8.505,2.042);
  \draw[{Stealth[length=3pt]}-{Stealth[length=3pt]},line width=0.7pt] (8.505,0.694) -- (8.505,2.028);
  \node[ecnode,anchor=west] at (8.6,1.35) {$\ch_{GR}$};
  \node[ecnode,below] at (8.505,0) {$\kh_{GR}$};
  \node[ecnode] (tg) at (4.6,2.5) {tangent, slope $\delta+n+x$};
  \draw[thin,black!55] (tg.south) -- (5.0,1.73);
\end{tikzpicture}
\caption{The Golden Rule ($\alpha = 1/3$, $\delta+n+x = 0.08$). Steady-state consumption is the gap between $\yh$ and the break-even line; it is largest where the slope of $\yh = \kh^{\alpha}$ equals $\delta+n+x$. With $s = \alpha$ the investment curve $\alpha\kh^{\alpha}$ crosses the break-even line exactly at $\kh_{GR}$.}
\end{figure}

\begin{extra}[Beyond the slides --- the Golden Rule in words]
The first-order condition says $f'(\kh_{GR}) = \alpha\kh_{GR}^{\alpha-1} = \delta+n+x$, i.e.\ $r - \delta = n + x$: at the Golden Rule, the net marginal product of capital equals the growth rate of the economy. The second derivative $\alpha(\alpha-1)(\kh^*)^{\alpha-2} < 0$ confirms this is a maximum.
\end{extra}

\subsection{Saving above or below the Golden Rule}
\begin{center}
\renewcommand{\arraystretch}{1.3}
\begin{tabular}{@{}p{0.46\linewidth}@{\qquad}p{0.46\linewidth}@{}}
\toprule
$s > \alpha$: \textbf{overaccumulation} & $s < \alpha$: \textbf{underaccumulation} \\
\midrule
$\kh^* > \kh_{GR}$. Lowering $s$ raises consumption \emph{immediately and} in the new steady state. The economy is \textbf{dynamically inefficient}. &
$\kh^* < \kh_{GR}$. Raising $s$ lowers consumption \emph{initially} but raises it \emph{eventually}. Whether that improves welfare depends on how future consumption is valued. \\
\bottomrule
\end{tabular}
\end{center}
In both cases $\ch^*(\kh) = \kh^{\alpha} - \gam\kh$ reaches its maximum at $\kh_{GR}$.

\subsection{A higher saving rate: the transition}
At date $t_0$, $s$ rises permanently from $s_0$ to $s_1$. $\kh$ is predetermined (the capital stock cannot jump).
\begin{flatnum}
  \item Investment $s\yh$ rises at once, so $\dk$ becomes positive.
  \item $\kh$ and $\yh$ rise gradually toward their higher steady-state levels.
  \item Growth of $Y/L$ temporarily exceeds $x$, then returns to $x$.
  \item $\ch = (1-s)\yh$ falls on impact; its eventual level may rise or fall, depending on the initial and new saving rates.
\end{flatnum}
\begin{keyresult}[Level vs.\ growth]
A higher saving rate raises the \emph{level} of output per worker, but does not permanently raise its \emph{growth rate} in the Solow model.
\end{keyresult}

\begin{extra}[Beyond the slides --- size of the jump]
Starting from the old steady state, $s_0(\kh_0^*)^{\alpha-1} = \gam$, so right after the change
\[
\left.\gr{\kh}\right|_{t_0} = s_1(\kh_0^*)^{\alpha-1} - \gam = \left(\frac{s_1}{s_0} - 1\right)\gam,
\qquad
\left.\gr{y}\right|_{t_0} = x + \alpha\left(\frac{s_1}{s_0} - 1\right)\gam.
\]
With $s_0 = 0.2$, $s_1 = 0.3$, $\alpha = 1/3$, $\gam = 0.08$: growth of $Y/L$ jumps from $2\%$ to about $3.3\%$ and then decays back to $2\%$ (figure below). Since $s_1 = 0.3 < \alpha$, consumption ends up higher than before.
\end{extra}

\begin{figure}[ht]
\centering
\begin{tikzpicture}[
  declare function={lam=0.0533333; z(\t)=3.75-1.25*exp(-lam*\t);},
  ttl/.style={font=\footnotesize,align=center,anchor=north}]
  \begin{scope}[xscale=0.045]
    \draw[ecaxis] (-20,0) -- (68,0);
    \draw[ecaxis] (-20,0) -- (-20,2.9);
    \node[ecnode,anchor=west] at (68,0) {$t$};
    \node[ecnode,anchor=south] at (-20,2.9) {$\ln(Y/L)$};
    \draw[budget] (-20,{1.4*0.058}) -- (0,{1.4*0.458});
    \draw[budget] plot[domain=0:60,samples=60] (\x,{1.4*(0.02*\x+0.5*ln(z(\x)))});
    \draw[black!55,dashed] (0,{1.4*0.458}) -- (60,{1.4*1.658});
    \draw[densely dotted] (0,0) -- (0,{1.4*0.458});
    \node[ecnode,below] at (0,0) {$t_0$};
    \node[ttl] at (24,-0.45) {(a) level shifts up;\\ slope returns to $x$};
  \end{scope}
  \begin{scope}[xshift=5.3cm,xscale=0.045]
    \draw[ecaxis] (-20,0) -- (68,0);
    \draw[ecaxis] (-20,0) -- (-20,2.9);
    \node[ecnode,anchor=west] at (68,0) {$t$};
    \node[ecnode,anchor=south] at (-20,2.9) {$\dot{y}/y$ (\%)};
    \draw[budget] (-20,1.2) -- (0,1.2);
    \draw[budget] plot[domain=0:60,samples=60]
         (\x,{0.6*100*(0.02+0.5*lam*1.25*exp(-lam*\x)/z(\x))});
    \draw[densely dotted] (0,0) -- (0,2.0);
    \draw[black!55,dashed] (0,1.2) -- (60,1.2);
    \node[ecnode,left] at (-20,1.2) {$2$};
    \node[ecnode,left] at (-20,2.0) {$3.3$};
    \draw[thin] (-21.5,2.0) -- (-18.5,2.0);
    \node[ecnode,below] at (0,0) {$t_0$};
    \node[ttl] at (24,-0.45) {(b) growth jumps,\\ decays back to $x$};
  \end{scope}
  \begin{scope}[xshift=10.6cm,xscale=0.045]
    \draw[ecaxis] (-20,0) -- (68,0);
    \draw[ecaxis] (-20,0) -- (-20,2.9);
    \node[ecnode,anchor=west] at (68,0) {$t$};
    \node[ecnode,anchor=south] at (-20,2.9) {$\ch$};
    \draw[budget] (-20,{6*0.265}) -- (0,{6*0.265});
    \draw[budget] plot[domain=0:60,samples=60] (\x,{6*(0.7*sqrt(z(\x))-1)});
    \draw[densely dotted] (0,0) -- (0,{6*0.265});
    \draw[black!55,dashed] (0,{6*0.356}) -- (60,{6*0.356});
    \node[ecnode,below] at (0,0) {$t_0$};
    \node[ttl] at (24,-0.45) {(c) falls on impact;\\ here ends higher ($s_1 < \alpha$)};
  \end{scope}
\end{tikzpicture}
\caption{Permanent increase in $s$ from $0.20$ to $0.30$ at date $t_0$ ($\alpha=1/3$, $\delta+n+x=0.08$, $x=0.02$). \emph{(Illustration, not on the slides.)}}
\end{figure}

\subsection{Foreign aid}
\topic{Consumption goods}
A temporary transfer raises consumption while it lasts. If it leaves saving, capital and productivity unchanged, it leaves the economy's \textbf{production path unchanged}.

\topic{Capital goods or investment finance}
If aid finances productive investment while $\kh$ is below its steady state, it speeds up capital accumulation during the transfer. After temporary aid ends, capital converges back toward the \textbf{original} steady state --- unless the transfer changes domestic saving or productivity persistently.

\topic{How a temporary investment transfer works}
Let $a_t$ be aid per effective worker, with a fraction $\varphi$ invested directly. If the remaining share $(1-\varphi)$ is treated as income and saved at rate $s$:
\[
\dk_t = s\kh_t^{\alpha} + \underbrace{[\varphi + s(1-\varphi)]\,a_t}_{\text{investment financed by aid}} - \gam\kh_t.
\]
\begin{flatlist}
  \item Cash aid entirely treated as disposable income: $\varphi = 0$, extra investment $sa_t$.
  \item In-kind investment aid: $\varphi = 1$, extra investment $a_t$.
  \item If $a_t = 0$ after a finite period, the equation returns to its original form, and the lasting effect on $\kh$ vanishes as the economy approaches its original steady state.
\end{flatlist}

\subsection{Technical assistance and the level of technology}
If assistance permanently raises labour-augmenting productivity $A_t$, output per worker $Y_t/L_t = A_tf\big(K_t/(A_tL_t)\big)$ can be permanently higher. At the instant $A_t$ rises, physical capital $K_t$ is fixed, so $\kh_t = K_t/(A_tL_t)$ \emph{falls}. Capital then accumulates back toward the original steady-state value of $\kh$ (if $s, n, x, \delta$ are unchanged).

\begin{keyresult}[Level vs.\ growth, again]
A permanent one-time increase in the \emph{level} of $A_t$ raises the long-run level of output per worker. A permanently higher \emph{growth rate} requires a persistent increase in the growth rate $x$ itself.
\end{keyresult}

\begin{extra}[Beyond the slides --- a worked version]
Let $A$ rise to $\mu A$ ($\mu > 1$) at $t_0$, starting from $\kh^*$. Then $\kh$ drops to $\kh^*/\mu$ and output per worker moves as
\[
A(\kh^*)^{\alpha} \;\xrightarrow{\text{impact}}\; \mu A\left(\frac{\kh^*}{\mu}\right)^{\!\alpha} = \mu^{1-\alpha}A(\kh^*)^{\alpha} \;\xrightarrow{\text{long run}}\; \mu A(\kh^*)^{\alpha}:
\]
an immediate jump by the factor $\mu^{1-\alpha}$, then capital accumulation takes it the rest of the way to $\mu$ times the original level.
\end{extra}

% =========================================================================
\section{Overall Assessment and Limitations}
\slides{L1 14--15; EV 18}

\subsection{What the model gets right}
Do the variables in the model behave like those in the data? The model is consistent with several long-run facts (L1 slide 14):
\begin{flatnum}
  \item Output per capita grows at a constant rate in the long run.
  \item $C/Y$ and $I/Y$ are constant in the long run.
  \item Capital grows at the constant rate $x$ faster than population in the long run.
  \item The rate of return to capital has no long-run trend.
  \item It can generate large and persistent differences in income levels and growth rates through differences in $s$, $A$ and $x$.
  \item[6--7.] It predicts conditional, but not absolute, convergence.
\end{flatnum}

\subsection{Overall assessment}
\begin{flatlist}
  \item It reproduces key balanced-growth facts for developed economies.
  \item It predicts conditional convergence through diminishing returns to capital.
  \item Capital accumulation can generate temporarily high growth, but not permanently high growth.
  \item Plausible saving and population-growth differences explain only part of international income gaps.
  \item Technology can account for large income differences, but the model leaves technology exogenous.
\end{flatlist}

\subsection{Limitations of the baseline model}
\begin{flatnum}
  \item \textbf{Key components are exogenous.} The saving rate, productivity and technological change are all exogenous, so the explanation is incomplete, and the model is not well suited to policy analysis.
  \item \textbf{Expenditure ratios.} $C/Y$ and $I/Y$ are also constant during the transitional growth period, which is counterfactual.
  \item \textbf{Convergence is too fast.} The implied speed of convergence is too high ($\beta \approx 5.3\%$, half-life $\approx 13$ years, vs.\ roughly $2\%$ and $35$ years in the data), motivating extensions of the basic framework.
\end{flatnum}

\begin{keyresult}[Central lesson]
The Solow model organizes the \emph{proximate} sources of growth and income differences (capital accumulation and technology), while leaving their deeper causes outside the model.
\end{keyresult}

% =========================================================================
\newpage
\section*{Formula Sheet}
\addcontentsline{toc}{section}{Formula Sheet}
\markboth{Formula Sheet}{}

{\renewcommand{\arraystretch}{1.55}
\begin{tabular}{@{}p{0.34\linewidth}p{0.62\linewidth}@{}}
\toprule
Growth-rate rules & $\dfrac{(\dot{XY})}{XY} = \gr{X} + \gr{Y}$, \quad $\dfrac{(\dot{X/Y})}{X/Y} = \gr{X} - \gr{Y}$, \quad $\dfrac{(\dot{X^a})}{X^a} = a\gr{X}$ \\
Exogenous growth & $L_t = L_0e^{nt}$, \quad $A_t = A_0e^{xt}$, \quad $A_tL_t$ grows at $n + x$ \\
Production & $Y = (AL)^{1-\alpha}K^{\alpha}$, \quad $y = A^{1-\alpha}k^{\alpha}$, \quad $\yh = \kh^{\alpha}$ \\
Factor prices & $r = \alpha\kh^{\alpha-1}$, \quad $w = (1-\alpha)A\kh^{\alpha}$, \quad shares $\alpha$ and $1-\alpha$ \\
Consumption & $\ch = (1-s)\yh$ \\
Solow equation & $\dk = s\kh^{\alpha} - \gam\kh$ \\
Steady state & $\kh^* = \left(\dfrac{s}{\delta+n+x}\right)^{\frac{1}{1-\alpha}}$, \quad $y^* = A\left(\dfrac{s}{\delta+n+x}\right)^{\frac{\alpha}{1-\alpha}}$ \\
Growth decomposition & $\gr{y} = x + \alpha\gr{\kh}$, \quad $\gr{\kh} = s\kh^{\alpha-1} - \gam$ \\
Long run & per capita grows at $x$; aggregates at $n + x$; $r$, $K/Y$, $C/Y$, $I/Y$ constant \\
Linear ODE & $\dot z + az = b \;\Rightarrow\; z_t = z^* + (z_0 - z^*)e^{-at}$, \quad $z^* = b/a$ \\
Exact Solow path & $\kh_t = \left[z^* + (\kh_0^{1-\alpha} - z^*)e^{-\beta t}\right]^{\frac{1}{1-\alpha}}$, \quad $z^* = \dfrac{s}{\delta+n+x}$ \\
Speed of convergence & $\beta = (1-\alpha)(\delta+n+x)$ \\
Log-linear path & $\tilde k_t - \tilde k_0 \simeq (1 - e^{-\beta t})(\tilde k^* - \tilde k_0)$, \quad $\tilde k = \ln\kh$ \\
Income convergence & $\ln y_t - \ln y_0 \simeq xt + (1 - e^{-\beta t})\left[\ln\yh^* - (\ln y_0 - \ln A_0)\right]$ \\
Growth regression & $\dfrac{\ln y_{i,T} - \ln y_{i,0}}{T} \simeq c_i - \lambda_T\ln y_{i,0} + \lambda_T\tfrac{\alpha}{1-\alpha}\ln s_i - \lambda_T\tfrac{\alpha}{1-\alpha}\ln(n_i + x + \delta)$, \quad $\lambda_T = \tfrac{1 - e^{-\beta T}}{T}$ \\
Half-life & $T_{1/2} = \ln 2/\beta$ \quad ($\approx 13$ years at $\beta = 0.0533$; $\approx 35$ at $0.02$) \\
Golden Rule & $\kh_{GR} = \left(\dfrac{\alpha}{\delta+n+x}\right)^{\frac{1}{1-\alpha}}$, \quad $s_{GR} = \alpha$ \\
Investment aid & $\dk = s\kh^{\alpha} + [\varphi + s(1-\varphi)]a - \gam\kh$ \\
\bottomrule
\end{tabular}}

\begin{transcriptnotes}
  \item \textbf{Sources}: \texttt{solow\_growth\_model\_slides.pdf} (L1, 15 slides), \texttt{solow\_growth\_model\_evaluation.pdf} (EV, 18), \texttt{solow\_differential\_equations.pdf} (DE, 14), \texttt{solow\_policy\_slides.pdf} (PO, 8).
  \item \textbf{2026-09-30 readability rewrite}: sections reordered to follow the decks (L1 $\to$ EV $\to$ DE $\to$ PO); every derivation step shown with its reason; key results boxed; material not on the slides moved into dashed ``Beyond the slides'' boxes; formula sheet added.
  \item \textbf{Corrections vs.\ previous version}: in the growth-regression derivation the first term is $xT/T = x$ (was printed as $xt/T$); a stray Markdown \texttt{**bold**} fixed; ``Kaldor facts'' label and ``inputs are essential'' removed (not on the slides) in favour of the slides' own list and wording; limitations now follow L1 slide 15 (three items).
  \item \textbf{Slide typos not copied}: EV slide 6 writes $\alpha\hat K^{\alpha-1}$ (should be $\alpha\kh^{\alpha-1}$); L1 slide 15 says ``contacts'' (``constant'').
\end{transcriptnotes}

\end{document}
